\documentclass[11pt,a4paper]{article}
\usepackage[margin=22mm]{geometry}
\usepackage[T1]{fontenc}
\usepackage{lmodern}
\usepackage{microtype}
\usepackage{xcolor}
\usepackage{booktabs}
\usepackage{tabularx}
\usepackage{longtable}
\usepackage{array}
\usepackage{amsmath}
\usepackage{enumitem}
\usepackage{listings}
\usepackage{fancyhdr}
\usepackage{hyperref}
\usepackage{graphicx}
\usepackage[section]{placeins}
\usepackage{float}

\definecolor{navy}{HTML}{183B56}
\definecolor{teal}{HTML}{087F83}
\definecolor{pale}{HTML}{EDF5F7}
\definecolor{codegray}{HTML}{F4F6F8}
\hypersetup{colorlinks=true,linkcolor=navy,urlcolor=teal,citecolor=teal,pdftitle={Client-Centric Consistency in MongoDB}}
\setlength{\parindent}{0pt}
\setlength{\parskip}{5pt}
\setlength{\emergencystretch}{2em}
\setlist{nosep,leftmargin=5mm}
\newcolumntype{Y}{>{\raggedright\arraybackslash}X}
\newcommand{\ok}{\textcolor{teal}{G}}
\newcommand{\noguarantee}{\textcolor{red!70!black}{NG}}
\newcommand{\na}{--}
\newcommand{\groupmembers}{\textit{Add group members and student IDs before submission}}

\pagestyle{fancy}
\fancyhf{}
\fancyhead[L]{\small MongoDB client-centric consistency lab}
\fancyhead[R]{\small DSA5208 project}
\fancyfoot[C]{\thepage}
\renewcommand{\headrulewidth}{0.3pt}

\lstset{basicstyle=\ttfamily\small,backgroundcolor=\color{codegray},frame=single,
  rulecolor=\color{black!15},breaklines=true,columns=fullflexible,showstringspaces=false}

\begin{document}

\begin{titlepage}
  \color{navy}
  \vspace*{22mm}
  {\Huge\bfseries Client-Centric Consistency in a Replicated MongoDB Database\par}
  \vspace{7mm}
  {\Large Experiments under replica lag, node failure, and network partition\par}
  \vspace{18mm}
  \begin{tabular}{@{}ll@{}}
    Course & DSA5208 Distributed Systems\\
    Database & MongoDB Community 7.0.32\\
    Deployment & Three Docker containers, one replica set\\
    Experiment date & 22 September 2026\\
    Group & \groupmembers
  \end{tabular}
  \vfill
  \color{black}
  \textbf{Submission note.} The report is generated from an audited formal run containing
  640 client histories and 2720 attempted operations. The accompanying archive includes the
  raw histories, three server logs, metadata, checksums, and reproduction scripts.
\end{titlepage}

\begin{abstract}
This project evaluates four client-centric consistency properties---read-your-writes,
monotonic reads, monotonic writes, and writes-follow-reads---in a three-member MongoDB
replica set. Eight client configurations form a complete $2\times2\times2$ matrix over
causal history tracking, read concern, and
write concern. The same application history is exercised under seven equally important
operating scenarios: normal operation, deliberate replica lag, secondary failure, primary
failure, a 2+1 network partition, complete 1+1+1 isolation, and partition-induced rollback.
The formal run contains 640 histories and 2720 attempted operations. Normal operation and
single-node failures established the healthy and failover baselines. Replica lag and both
partition patterns exposed the difference between stale successful reads and operations
that wait or fail to preserve causal order. The rollback scenario exposed the durability
consequences of acknowledged $w:1$ writes. Across the full matrix, the observations agree
with the predicted client-centric guarantees and demonstrate that consistency and
availability must be interpreted together.
\end{abstract}

\tableofcontents
\newpage

\section{System and deployment}

\subsection{Chosen database}
MongoDB is a document database whose replica sets provide automatic primary election,
asynchronous replication, configurable read and write concerns, and causally consistent
client sessions. These mechanisms make it possible to vary consistency and durability on
a per-operation basis while preserving the same application workload \cite{mongo-causal,mongo-read,mongo-write}.

The experiment uses MongoDB Community 7.0.32 and PyMongo 4.10.1. All versions are pinned
in the Dockerfile and \texttt{requirements.txt}. The storage engine is WiredTiger.

\subsection{Architecture}

\begin{center}
\fcolorbox{teal}{pale}{\begin{minipage}{0.90\linewidth}
\centering
\textbf{Python experiment controller}\par
\smallskip
$\downarrow$ host ports 29101--29103\par
\smallskip
\begin{tabular}{ccc}
\textbf{mongo0} & \textbf{mongo1} & \textbf{mongo2}\\
electable, vote 1 & electable, vote 1 & priority 0, vote 1\\
own process/IP/volume & own process/IP/volume & own process/IP/volume
\end{tabular}\par
\smallskip
\textit{Replication and heartbeats use a private Docker bridge network.}
\end{minipage}}
\end{center}

Each container stores a complete data copy in a named volume. All three members vote, but
\texttt{mongo2} has priority zero so that it cannot become primary. Nodes 0 and 1 are
electable. The replica set therefore survives one node failure and can elect a primary on
the two-node side of a 2+1 split. The controller uses direct connections so that each read
target is known. Causal metadata is explicitly transferred when the logical client changes
physical connections; a PyMongo session is never passed to a different \texttt{MongoClient}.

This is a distributed logical deployment using multiple containers on one physical Mac,
as permitted by the assignment. It does not emulate independent power, disks, or host
kernels.

\subsection{Installation and startup}

Docker Desktop and Docker Compose are required. From the repository root:

\begin{lstlisting}[language=bash]
python3 -m venv .venv
.venv/bin/python -m pip install -r requirements.txt
.venv/bin/python -m unittest -v test_checker.py
.venv/bin/python run.py --work-dir work/matrix8-timeout5s-full-02 \
  --out results/matrix8-timeout5s-full-02
.venv/bin/python analyze.py results/matrix8-timeout5s-full-02
\end{lstlisting}

The runner builds the pinned MongoDB image, creates fresh containers and volumes, calls
the replica-set initialization command, waits for one primary and two secondaries, performs the tests,
heals all partitions, collects logs, and stops the containers. New directories are required
for every run, preventing accidental reuse of earlier data.

\section{Consistency configurations and predictions}

MongoDB read concern specifies which data a read may expose, while write concern specifies
how many voting members must acknowledge a write. A causally consistent session carries
logical time from an earlier operation to a later operation. For reads, the driver sends an
\texttt{afterClusterTime} lower bound so that a server cannot return an older snapshot merely
to respond quickly \cite{mongo-causal,pymongo-session}.

\begin{table}[ht]
\centering
\caption{Client configurations.}
\begin{tabular}{@{}clll@{}}
\toprule
ID & Causal tracking & Read concern & Write concern\\
\midrule
A & off & local & $w:1$\\
B & on & local & $w:1$\\
C & off & local & $w:\text{majority}$\\
D & on & local & $w:\text{majority}$\\
E & off & majority & $w:1$\\
F & on & majority & $w:1$\\
G & off & majority & $w:\text{majority}$\\
H & on & majority & $w:\text{majority}$\\
\bottomrule
\end{tabular}
\end{table}

The table is ordered as four adjacent matched pairs: A/B, C/D, E/F, and G/H. Within
each pair, read concern and write concern are fixed and only causal tracking is
changed. This makes the causal-session effect identifiable without changing the
other two client settings at the same time.

Table~\ref{tab:predictions} records the predictions made from the documented causal-session
matrix. G means a documented guarantee for successful operations under the stated settings;
NG means that a permitted execution can violate the property. An NG prediction does not
mean every execution must show a violation.

\begin{table}[ht]
\centering
\caption{Predicted client-centric guarantees.}
\label{tab:predictions}
\begin{tabular}{@{}lcccc@{}}
\toprule
Configuration & RYW & MR & MW & WFR\\
\midrule
A: local/$w:1$, causal off & \noguarantee & \noguarantee & \noguarantee & \noguarantee\\
B: local/$w:1$, causal on & \noguarantee & \noguarantee & \noguarantee & \noguarantee\\
C: local/majority, causal off & \noguarantee & \noguarantee & \noguarantee & \noguarantee\\
D: local/majority, causal on & \ok & \noguarantee & \ok & \noguarantee\\
E: majority/$w:1$, causal off & \noguarantee & \noguarantee & \noguarantee & \noguarantee\\
F: majority/$w:1$, causal on & \noguarantee & \ok & \noguarantee & \ok\\
G: majority/majority, causal off & \noguarantee & \noguarantee & \noguarantee & \noguarantee\\
H: majority/majority, causal on & \ok & \ok & \ok & \ok\\
\bottomrule
\end{tabular}
\end{table}

The report uses a durable-history interpretation: if a write returned an acknowledgment,
later loss of that write during rollback can witness an RYW or MW violation. MongoDB also
documents a weaker interpretation in which rollback is itself part of the history. We state
the chosen interpretation to avoid mixing the two definitions \cite{mongo-causal,mongo-rollback}.

\section{Experimental method}

\subsection{Workload and oracle}
Every trial uses a unique document initialized on all three members as
\texttt{\{a:0,b:0\}}. The usual operation order is:

\begin{center}
\texttt{W1: set a=1} $\rightarrow$ \texttt{R1} $\rightarrow$ \texttt{R2}
$\rightarrow$ \texttt{W2: set b=1 and return pre-image}.
\end{center}

\texttt{findOneAndUpdate(..., returnDocument=BEFORE)} makes W2 a real modifying write while
atomically revealing the database state on which it executed. This avoids incorrectly
inferring write ordering from a later final value.

\begin{table}[ht]
\centering
\caption{Operational checks for the four properties.}
\begin{tabularx}{\textwidth}{@{}p{31mm}YY@{}}
\toprule
Property & Eligible check & Violation witness\\
\midrule
Read-your-writes & W1 acknowledged and at least one later read returned & Any successful read invoked after W1 completed returns $a<1$\\
Monotonic reads & R1 and R2 both returned & $R2.a<R1.a$\\
Monotonic writes & W1 and W2 acknowledged & W2 pre-image has $a<1$\\
Writes-follow-reads & W2 acknowledged after a successful read & W2 pre-image is older than the maximum version observed by earlier reads\\
\bottomrule
\end{tabularx}
\end{table}

Read errors expose no value and therefore produce an inconclusive read-to-read check. A
write timeout can have an unknown outcome and is not treated as an acknowledged predecessor.
Errors and timeouts are reported separately from stale-value violations. Eight synthetic
unit tests validate preserved histories, stale reads, timeout handling, rollback, WFR order,
and the requirement that RYW examine every successful read following W1.

\subsection{Scenarios}

\begin{longtable}{@{}p{35mm}p{104mm}r@{}}
\caption{Operating scenarios and repetitions per configuration.}\\
\toprule
Scenario & Fault and routing & $n$\\
\midrule
\endfirsthead
\toprule Scenario & Fault and routing & $n$\\ \midrule
\endhead
Normal & No injected fault; R2 targets a known secondary. & 20\\
Replica lag & \texttt{mongo2} applies operations with a three-second delay; R2 targets it. & 10\\
Secondary crash & Kill \texttt{mongo2} with SIGKILL; the surviving pair remains available. & 10\\
Primary crash & W1/R1, kill primary, await election, then R2/W2 on the new primary. & 10\\
2+1 partition & Isolate the old primary; the other two nodes elect a new primary. W1/R1/W2 use the majority side and R2 probes the minority. & 10\\
1+1+1 isolation & Drop all peer traffic while host-to-container client access remains. & 10\\
Partition + rollback & Isolate before W1, kill old primary, heal survivors, elect a new primary, then execute W2/R2. & 10\\
\bottomrule
\end{longtable}

Faults are implemented with container-local \texttt{iptables} INPUT/OUTPUT rules using
\texttt{NET\_ADMIN}. Rules block only peer container addresses; the controller retains access
through published loopback ports. Every disruptive trial receives its own cut/heal cycle and
an administrative $w:3$ recovery barrier. The audit validates actual wire commands,
read/write concerns, causal bounds, unique trial IDs, and the expected topology of every
partition snapshot.

\section{Results}

The formal run completed 640 histories and 2720 attempted operations. Command monitoring
captured 2700 wire commands; 20 rollback operations were rejected by PyMongo before any
command was sent because a restarted direct connection temporarily reported that sessions
were unsupported. The audit records these separately and found no missing wire evidence for
any successful operation, no incorrect concern parameters, and no failure among 160 isolated
topology snapshots.

The tables use one common notation. RC and WC denote read and write concern; property cells
are \emph{violations/eligible checks}; and ``--'' means that a failed prerequisite made the
check inconclusive. Normal operation used 20 four-operation histories per configuration, the
other non-partition baselines used 10, and each network-partition history added a fifth
unavailable-write probe. Thus the error denominators are 80, 40, or 50 as appropriate.
Errors measure availability, not consistency violations. A recurring E/F contrast is also
useful: both use majority reads and $w:1$, but only F propagates causal time. E may therefore
return an older committed snapshot after acknowledging W1, whereas F must wait or fail rather
than return below its causal bound.

\subsection{Normal operation}
\begin{table}[H]
\centering\scriptsize
\caption{Normal operation: violations/eligible checks and errors/operations.}
\begin{tabular}{@{}lcllrrrrr@{}}
\toprule Config & Causal & RC & WC & RYW & MR & MW & WFR & Errors\\ \midrule
A & off & local & 1 & 0/20 & 0/20 & 0/20 & 0/20 & 0/80\\
B & on  & local & 1 & 0/20 & 0/20 & 0/20 & 0/20 & 0/80\\
C & off & local & majority & 0/20 & 0/20 & 0/20 & 0/20 & 0/80\\
D & on  & local & majority & 0/20 & 0/20 & 0/20 & 0/20 & 0/80\\
E & off & majority & 1 & 20/20 & 0/20 & 0/20 & 0/20 & 0/80\\
F & on  & majority & 1 & 0/20 & 0/20 & 0/20 & 0/20 & 0/80\\
G & off & majority & majority & 0/20 & 0/20 & 0/20 & 0/20 & 0/80\\
H & on  & majority & majority & 0/20 & 0/20 & 0/20 & 0/20 & 0/80\\ \bottomrule
\end{tabular}
\end{table}

All nodes were healthy, R2 targeted \texttt{mongo2}, and all 640 operations succeeded.
Only E violated a property: in every history, at least one majority read returned the
committed snapshot preceding its acknowledged but not-yet-majority-committed W1. Its two
reads did not move from new to old, so MR remained unviolated. F's causal lower bound removed
that stale-successful execution. This E/F comparison is direct evidence that majority read
concern alone does not provide RYW. The remaining zero cells mean only that this run observed
no witness, not that every possible execution is safe.
\FloatBarrier

\subsection{Replica lag}
\begin{table}[H]
\centering\scriptsize
\caption{Three-second replica lag: violations/eligible checks and errors/operations.}
\begin{tabular}{@{}lcllrrrrr@{}}
\toprule Config & Causal & RC & WC & RYW & MR & MW & WFR & Errors\\ \midrule
A & off & local & 1 & 10/10 & 10/10 & 0/10 & 0/10 & 0/40\\
B & on  & local & 1 & 0/10 & 0/10 & 0/10 & 0/10 & 0/40\\
C & off & local & majority & 10/10 & 10/10 & 0/10 & 0/10 & 0/40\\
D & on  & local & majority & 0/10 & 0/10 & 0/10 & 0/10 & 0/40\\
E & off & majority & 1 & 10/10 & 0/10 & 0/10 & 0/10 & 0/40\\
F & on  & majority & 1 & 0/10 & 0/10 & 0/10 & 0/10 & 0/40\\
G & off & majority & majority & 10/10 & 10/10 & 0/10 & 0/10 & 0/40\\
H & on  & majority & majority & 0/10 & 0/10 & 0/10 & 0/10 & 0/40\\ \bottomrule
\end{tabular}
\end{table}

\texttt{mongo2} applied oplog entries three seconds late and served R2. Non-causal A, C, E,
and G returned a value that omitted W1, giving 10/10 RYW violations. A, C, and G first read
the new value on the primary and then the old value on the delayed member, so they also gave
10/10 MR regressions. E instead read the same older majority snapshot twice: it failed RYW
but did not move backward. Causal B, D, F, and H sent \texttt{afterClusterTime}; because the
five-second read deadline exceeded the finite three-second lag, they waited for the member to
catch up and then succeeded. All 320 operations completed, showing consistency obtained here
through added latency rather than reduced availability. MW and WFR remained preserved because
the scenario delayed replication but introduced neither rollback nor write reordering.
\FloatBarrier

\subsection{Secondary crash}
\begin{table}[H]
\centering\scriptsize
\caption{Secondary crash: violations/eligible checks and errors/operations.}
\begin{tabular}{@{}lcllrrrrr@{}}
\toprule Config & Causal & RC & WC & RYW & MR & MW & WFR & Errors\\ \midrule
A & off & local & 1 & 0/10 & 0/10 & 0/10 & 0/10 & 0/40\\
B & on  & local & 1 & 0/10 & 0/10 & 0/10 & 0/10 & 0/40\\
C & off & local & majority & 0/10 & 0/10 & 0/10 & 0/10 & 0/40\\
D & on  & local & majority & 0/10 & 0/10 & 0/10 & 0/10 & 0/40\\
E & off & majority & 1 & 10/10 & 0/10 & 0/10 & 0/10 & 0/40\\
F & on  & majority & 1 & 0/10 & 0/10 & 0/10 & 0/10 & 0/40\\
G & off & majority & majority & 0/10 & 0/10 & 0/10 & 0/10 & 0/40\\
H & on  & majority & majority & 0/10 & 0/10 & 0/10 & 0/10 & 0/40\\ \bottomrule
\end{tabular}
\end{table}

\texttt{mongo2} was killed before the histories began. The primary and the surviving
secondary still held two of three votes, so majority writes remained available and R2 could
use the surviving secondary. All 320 operations succeeded. E again produced 10/10 RYW
witnesses without MR regression through the shared $w:1$/majority-snapshot mechanism; F's
causal bound waited until that write was visible. No rollback or competing branch occurred,
so MW and WFR also remained unviolated. The key result is that losing one secondary did not
remove quorum, although it did not repair E's weak acknowledgment/read combination.
\FloatBarrier

\subsection{Primary crash and election}
\begin{table}[H]
\centering\scriptsize
\caption{Primary crash and election: violations/eligible checks and errors/operations.}
\begin{tabular}{@{}lcllrrrrr@{}}
\toprule Config & Causal & RC & WC & RYW & MR & MW & WFR & Errors\\ \midrule
A & off & local & 1 & 0/10 & 0/10 & 0/10 & 0/10 & 0/40\\
B & on  & local & 1 & 0/10 & 0/10 & 0/10 & 0/10 & 0/40\\
C & off & local & majority & 0/10 & 0/10 & 0/10 & 0/10 & 0/40\\
D & on  & local & majority & 0/10 & 0/10 & 0/10 & 0/10 & 0/40\\
E & off & majority & 1 & 10/10 & 0/10 & 0/10 & 0/10 & 0/40\\
F & on  & majority & 1 & 0/10 & 0/10 & 0/10 & 0/10 & 0/40\\
G & off & majority & majority & 0/10 & 0/10 & 0/10 & 0/10 & 0/40\\
H & on  & majority & majority & 0/10 & 0/10 & 0/10 & 0/10 & 0/40\\ \bottomrule
\end{tabular}
\end{table}

W1 and R1 ran on the original primary; after SIGKILL, the harness waited for a new primary
before issuing R2 and W2. All 320 measured operations therefore succeeded. This zero-error
result describes post-election recovery, not the duration or errors of the election window.
E again gave 10/10 RYW witnesses through the $w:1$/majority-snapshot gap; F and the other
configurations returned no contradictory successful value. Replication was healthy before
the crash, so this scenario demonstrates failover but does not prove that every $w:1$ write
survives. The rollback experiment isolates that stronger failure.
\FloatBarrier

\subsection{2+1 network partition}
\begin{table}[H]
\centering\scriptsize
\caption{2+1 partition: violations/eligible checks and errors/operations.}
\begin{tabular}{@{}lcllrrrrr@{}}
\toprule Config & Causal & RC & WC & RYW & MR & MW & WFR & Errors\\ \midrule
A & off & local & 1 & 10/10 & 10/10 & 0/10 & 0/10 & 10/50\\
B & on  & local & 1 & 0/10 & -- & 0/10 & 0/10 & 20/50\\
C & off & local & majority & 10/10 & 10/10 & 0/10 & 0/10 & 10/50\\
D & on  & local & majority & 0/10 & -- & 0/10 & 0/10 & 20/50\\
E & off & majority & 1 & 10/10 & 0/10 & 0/10 & 0/10 & 10/50\\
F & on  & majority & 1 & 0/10 & -- & 0/10 & 0/10 & 20/50\\
G & off & majority & majority & 10/10 & 10/10 & 0/10 & 0/10 & 10/50\\
H & on  & majority & majority & 0/10 & -- & 0/10 & 0/10 & 20/50\\ \bottomrule
\end{tabular}
\end{table}

The isolated old primary became the one-node minority, while the other two members elected a
new primary. W1, R1, and W2 succeeded on that majority side; R2 deliberately probed the stale
minority. Non-causal A, C, and G moved from R1's new value to R2's old value and violated both
RYW and MR. E again stayed on one older majority snapshot, violating only RYW. Causal B, D,
F, and H could not satisfy R2's majority-side \texttt{afterClusterTime} on the isolated node,
so R2 timed out and MR was inconclusive instead of stale.

The error totals follow from two probes: every history's minority-side write was rejected
with \texttt{NotPrimaryError} (80 errors), and the 40 causal R2 reads reached
\texttt{ExecutionTimeout}. Hence non-causal rows contain 10/50 errors and causal rows 20/50.
MW and WFR stayed at zero because their successful prerequisites ran in order on the
majority side. Unlike finite lag, a continuing partition could not clear within the deadline.
\FloatBarrier

\subsection{1+1+1 complete isolation}
\begin{table}[H]
\centering\scriptsize
\caption{1+1+1 isolation: violations/eligible checks and errors/operations.}
\begin{tabular}{@{}lcllrrrrr@{}}
\toprule Config & Causal & RC & WC & RYW & MR & MW & WFR & Errors\\ \midrule
A & off & local & 1 & 10/10 & 10/10 & 0/10 & 0/10 & 10/50\\
B & on  & local & 1 & 0/10 & -- & 0/10 & 0/10 & 20/50\\
C & off & local & majority & -- & 10/10 & -- & -- & 30/50\\
D & on  & local & majority & -- & -- & -- & -- & 40/50\\
E & off & majority & 1 & 10/10 & 0/10 & 0/10 & 0/10 & 10/50\\
F & on  & majority & 1 & -- & -- & 0/10 & -- & 30/50\\
G & off & majority & majority & -- & 0/10 & -- & -- & 30/50\\
H & on  & majority & majority & -- & -- & -- & -- & 50/50\\ \bottomrule
\end{tabular}
\end{table}

All peer paths were blocked, leaving three one-node components and no possible quorum. During
the old primary's brief step-down window, weak $w:1$ configurations remained partly
available: A observed a new-to-old local-read regression, while E's majority reads stayed on
the older snapshot. C's locally applied but unacknowledged W1 was visible to R1, so its later
old R2 value witnessed MR even though RYW and MW were ineligible. Stronger requirements
progressively replaced successful stale observations with failures; H returned no successful
set of prerequisites and therefore has dashes rather than consistency passes.

The row totals consist of 40 W1 and 30 W2 majority-write timeouts, 10 H W2
\texttt{NotPrimaryError}s, 20 causal R1 and 40 causal R2 timeouts, and 80 rejected
post-step-down probes. These counts explain the progression from 10/50 errors for A and E to
50/50 for H. The result is the intended safety--availability contrast: operations that could
not meet quorum or causal requirements failed instead of returning a contradictory value.
\FloatBarrier

\subsection{Partition followed by rollback}
\begin{table}[H]
\centering\scriptsize
\caption{Partition followed by rollback: violations/eligible checks and errors/operations.}
\begin{tabular}{@{}lcllrrrrr@{}}
\toprule Config & Causal & RC & WC & RYW & MR & MW & WFR & Errors\\ \midrule
A & off & local & 1 & 10/10 & 10/10 & 10/10 & 10/10 & 0/40\\
B & on  & local & 1 & 6/7 & 6/6 & 6/6 & 6/6 & 8/40\\
C & off & local & majority & -- & 9/9 & -- & 7/7 & 14/40\\
D & on  & local & majority & -- & 9/9 & -- & 3/3 & 18/40\\
E & off & majority & 1 & 9/9 & 0/8 & 8/8 & 0/8 & 4/40\\
F & on  & majority & 1 & 7/7 & -- & 9/9 & -- & 14/40\\
G & off & majority & majority & -- & 0/10 & -- & 0/6 & 14/40\\
H & on  & majority & majority & -- & -- & -- & -- & 34/40\\ \bottomrule
\end{tabular}
\end{table}

The nodes were isolated before W1. A $w:1$ write could be acknowledged only on the old
primary, which was then killed; the two survivors elected a branch without W1. W2's pre-image
and R2 exposed that loss. A completed this sequence in every history and violated all four
properties. B, E, and F produced the same rollback witnesses whenever their prerequisites
returned: causal metadata could constrain later operations, but it could not recreate a
write absent from the elected branch. E's majority reads stayed on the older snapshot, so its
MR and WFR checks remained unviolated even while RYW and MW failed.

Majority-write C, D, G, and H never acknowledged isolated W1; RYW and MW were therefore
inconclusive rather than violated. C and D could still read a locally applied W1 before the
timeout and later observe its absence, explaining their MR and eligible WFR witnesses. G's
successful majority reads stayed on the old snapshot, while H returned too few successful
prerequisites for any property check.

The 106 errors comprise 40 W1 and 18 W2 \texttt{WTimeoutError}s, 28 causal-read timeouts, and
20 PyMongo \texttt{ConfigurationError}s raised before transmission while restarted direct
connections temporarily lacked session capability. The latter explain the unequal eligible
denominators and are treated as a harness limitation. The central result is nevertheless
clear: causal tracking cannot restore a rolled-back $w:1$ predecessor, while majority write
concern avoids acknowledging that unsafe history at the cost of availability.
\FloatBarrier

\section{Comparison with predictions}

The observations agree with the documented boundaries and expose the purpose of each matched
pair:
\begin{itemize}
  \item A/B and C/D show that causal tracking prevents successful backward reads even with
  local read concern, although rollback can still invalidate a non-durable $w:1$ predecessor.
  \item E/F isolates causal tracking with majority reads and $w:1$. E returned stale committed
  snapshots; F waited or failed rather than cross its causal lower bound.
  \item G/H shows that majority concerns without causal tracking are not a substitute for a
  causal session. G violated RYW/MR under lag and 2+1 partition; H did not return a
  contradictory successful value.
  \item Comparing A/B/E/F ($w:1$) with C/D/G/H (majority write) shows why acknowledgment
  durability matters during rollback. Majority-write histories became inconclusive when
  quorum was absent instead of accepting the isolated predecessor.
\end{itemize}

Zero observed violations are not proofs of universal guarantees. Guarantees come from the
documented model; executions supply corroborating examples or counterexamples. Likewise,
timeouts are evidence about availability rather than successful stale observations.

\section{Limitations}

\begin{enumerate}
  \item All containers share one Mac and Docker's Linux VM, so host, kernel, disk, and power
  failures are correlated.
  \item Direct connections intentionally force particular read targets. The experiment does
  not characterize normal driver load balancing.
  \item A three-second delayed member is a laboratory control, not a realistic WAN model.
  \item Each trial uses one document, a binary monotone version field, and sequential client
  operations. Cross-document causality, sharding, transactions, and concurrent overwrites
  are outside scope.
  \item Fault scenarios use ten deterministic repetitions per configuration. This exposes
  timing variation but is still too small for production-frequency estimates.
  \item Repeated restart cycles caused 20 PyMongo pre-command session-capability errors in
  rollback. They are reported as availability failures, but they reduced some eligible
  denominators and are a harness limitation.
  \item The four property checks overlap within a history and must not be summed as
  statistically independent samples.
  \item The durable-history treatment of rollback is explicit but not the only terminology
  used in MongoDB documentation.
\end{enumerate}

\section{Reproducibility and submitted evidence}

The repository contains \texttt{compose.yaml}, \texttt{Dockerfile}, \texttt{lab.py},
\texttt{docker\_backend.py}, \texttt{analyze.py}, checker unit tests, dependency pins, and a
README. The formal A--H run is stored in this archive:
\begin{center}
\small\path{matrix8-timeout5s-full-20260922-raw.zip}
\end{center}

\begin{itemize}
  \item \texttt{history.jsonl}: operations, wire commands, returned values, causal tokens,
  fault events, topology snapshots, and recovery observations;
  \item \texttt{trials.json}: structured histories and four-property classifications;
  \item three \texttt{mongod-*.log} files containing elections, replication, and rollback;
  \item \texttt{metadata.json}, \texttt{audit.json}, and derived summaries.
\end{itemize}

\texttt{RAW\_DATA\_MANIFEST.json} supplies SHA-256 digests for every original file and the
ZIP. Re-analysis must report 640 histories, 2720 attempted operations, 2700 wire commands,
20 pre-command errors, 160 isolation snapshots, and zero parameter, wire-coverage, or
isolation audit failures.

\section{AI usage disclosure}

OpenAI Codex assisted with interpreting the assignment and MongoDB documentation, designing
the experiment matrix, implementing and debugging the harness, executing local Docker tests,
auditing raw evidence, and drafting this report. Database measurements were produced by real
MongoDB processes. The group remains responsible for reviewing the code, validating the
results, adding member information, and submitting the final work.

\begin{thebibliography}{9}
\bibitem{mongo-causal} MongoDB, ``Causal Consistency and Read and Write Concerns,'' MongoDB 7.0 Manual. \url{https://www.mongodb.com/docs/v7.0/core/causal-consistency-read-write-concerns/}
\bibitem{mongo-read} MongoDB, ``Read Isolation, Consistency, and Recency,'' MongoDB 7.0 Manual. \url{https://www.mongodb.com/docs/v7.0/core/read-isolation-consistency-recency/}
\bibitem{mongo-write} MongoDB, ``Write Concern,'' MongoDB 7.0 Manual. \url{https://www.mongodb.com/docs/v7.0/reference/write-concern/}
\bibitem{mongo-deploy} MongoDB, ``Deploy a Self-Managed Replica Set for Testing and Development,'' MongoDB 7.0 Manual. \url{https://www.mongodb.com/docs/v7.0/tutorial/deploy-replica-set-for-testing/}
\bibitem{mongo-delay} MongoDB, ``Delayed Replica Set Members,'' MongoDB 7.0 Manual. \url{https://www.mongodb.com/docs/v7.0/core/replica-set-delayed-member/}
\bibitem{mongo-rollback} MongoDB, ``Rollbacks During Replica Set Failover,'' MongoDB 7.0 Manual. \url{https://www.mongodb.com/docs/v7.0/core/replica-set-rollbacks/}
\bibitem{pymongo-session} MongoDB, ``PyMongo ClientSession API,'' PyMongo 4.10.1. \url{https://pymongo.readthedocs.io/en/4.10.1/api/pymongo/client_session.html}
\bibitem{docker-compose} Docker, ``Compose file reference: services.'' \url{https://docs.docker.com/reference/compose-file/services/}
\bibitem{lecture} DSA5208, ``Lecture 3: Consistency Models,'' course lecture slides, 2026.
\end{thebibliography}

\appendix
\section{Exact configuration summary}

MongoDB uses replica set \texttt{cc\_lab}, a 128 MB oplog, and a 0.25 GB cache per node.
Election/heartbeat intervals are 15000/500 ms. Read, write-concern, and socket deadlines are
5000/5000/7000 ms; retries are disabled and the random seed is 419. Recovery barriers use
longer administrative timeouts than measured operations.

\end{document}
